\documentclass[10pt,a4paper]{amsart}
\usepackage[margin=19mm]{geometry}
\usepackage{amsmath,amssymb,mathtools}
\usepackage[scr=rsfs,cal=euler]{mathalfa}
\usepackage{xfrac}
\usepackage[unicode,hidelinks,bookmarks=false]{hyperref}
\newcommand{\ie}{i.e.\ }
\newcommand{\cf}{cf.\ }
\newcommand{\resp}{respectively\ }
\newcommand{\Subsection}{\subsection}
\providecommand{\nobreakdash}{\nobreak\hbox{-}}
\setlength{\emergencystretch}{2em}
\multlinegap0pt
\usepackage{amssymb}
\usepackage{mathtools}
\usepackage[shortlabels]{enumitem}
\newtheorem{lemma}{Lemma}
\newtheorem{definition}{Definition}
\newcommand{\Pp}{\mathbb P}
\newcommand{\PsiThree}{\Psi_{3}}
\title{Martin's Axiom, Large Continuum and Global $\Sigma^1_n$-Uniformization}
\author{Stefan Hoffelner}
\date{Source: arXiv:2605.21189v1 (CC BY 4.0)}
\begin{document}
\maketitle

\noindent\emph{Source and licence.} Martin's Axiom, Large Continuum and Global $\Sigma^1_n$-Uniformization. Version arXiv:2605.21189v1 (CC BY 4.0), distributed by arXiv under the Creative Commons Attribution 4.0 International licence, http://creativecommons.org/licenses/by/4.0/, as recorded in the arXiv licence metadata for this exact version and shown on the abstract page https://arxiv.org/abs/2605.21189v1. Full licence notice: \texttt{licenses/arxiv-2605.21189-CC-BY-4.0.txt}.\par\smallskip

\noindent\textit{Editorial scope.} Counted unit: the article's lemma lem:block-correctness (source0.tex lines 701--709). The article's complete original proof of it is delivered with it (source0.tex lines 711--846, one \texttt{proof} environment of 136 lines). No mathematics is written, repaired or re-derived: the statement and the proof are the article's own lines, in the article's own notation, and no retained line is substituted or re-flowed.  Retained uncounted as the source-local context the proof needs: lines 599--609; lines 665--681.  Nothing was omitted from any retained span, so the package carries no omission record.  No retained line contains a citation, so the wrapper carries no bibliography and no bibliography entry.  No retained line contains a figure environment, an image-inclusion command or drawing code, so no image is delivered and the package declares no asset dependency.  Editorial formatting only, in addition: an AMS \texttt{amsart} preamble that restates the article's own declarations and macros used by the retained lines, namely the environments \texttt{lemma}, \texttt{definition} on separate counters, together with the standard packages those lines need.  The retained environments print in this document as Lemma 1, Definition 1 in order of appearance, whereas the article prints the same items with the numbering of its own full document.  The complete native source of the article as distributed in the arXiv e-print for this exact version is bundled unchanged as \texttt{sources/201050/source0.tex}.
\begin{lemma}\label{lem:no-accidental-codes}
Let \(\alpha\in\mathcal L_{\mathrm{WO}}\cup\mathcal L_\Sigma\) be a fresh block,
and let \(\mathbb R\) be a c.c.c. side forcing supported below \(\alpha\) in the
above sense.  Then, after forcing with \(\mathbb R\), the fresh block
\[
   \langle T_{\omega\cdot\alpha+n}:n<\omega\rangle
\]
which is added by the tail preliminary factor still consists of Suslin trees.
Consequently no branch-versus-specialize pattern appears on the block
\(\alpha\) before the construction deliberately acts on that block.
\end{lemma}

\begin{definition}[The base predicate]\label{def:psi3}
For a real $z$, let $\PsiThree(z)$ be the assertion that there is a real $R$ such
that for every countable suitable model $M$ with $R,z\in M$, there is
$\bar\alpha\in\mathcal L_\Sigma^M$ such that, for every $n\in\omega$, the tree
\[
T^M_{\bar\alpha,n}=\{\gamma<\omega_1^M:
S^M_{\omega_1\cdot(\omega\cdot\bar\alpha+n)+\gamma}
\text{ is nonstationary in }M\}
\]
is an $\omega_1^M$-tree and $M$ sees that
\[
\begin{array}{rcl}
n\in\Delta(z)&\Longrightarrow& T^M_{\bar\alpha,n}\text{ is specialized},\\
n\notin\Delta(z)&\Longrightarrow& T^M_{\bar\alpha,n}\text{ has a cofinal branch}.
\end{array}
\]
\end{definition}

\begin{lemma}[Unary exactness of the base code]\label{lem:block-correctness}
Let $\Pp$ be a completed iteration satisfying this block discipline, and let $G$ be
$\Pp$-generic.  Then, in $V[G]$, for every real $z$,
\[
\PsiThree(z)
\]
holds if and only if $z$ was deliberately placed by the iteration into one of the
uniformization-reserved branch-versus-specialize coding blocks.
\end{lemma}

\begin{proof}
We first fix the notation which will be used in both directions.  If
$\alpha\in\mathcal L_{\Sigma}$ and $n\in\omega$, put
\[
E^G_{\alpha,n}=\{\gamma<\omega_1:
S_{\omega_1\cdot(\omega\cdot\alpha+n)+\gamma}\text{ is nonstationary in }V[G]\}.
\]
By the preliminary FFDZ coding and the preservation theorem for the stationary sets,
\[
E^G_{\alpha,n}=T_{\omega\cdot\alpha+n}.                 \tag{1}
\]
More precisely, if $\gamma\in T_{\omega\cdot\alpha+n}$, the preliminary stationary-killing
forcing added a club through the complement of
$S_{\omega_1\cdot(\omega\cdot\alpha+n)+\gamma}$; while if
$\gamma\notin T_{\omega\cdot\alpha+n}$, the corresponding stationary set remains
stationary through the preliminary forcing and through all later c.c.c. stages which do
not explicitly kill it.  Thus the nonstationarity pattern recovers exactly the tree
$T_{\omega\cdot\alpha+n}$.

For a uniformization block $\alpha$, define its final pattern $b_\alpha\subseteq\omega$ by
\[
n\in b_\alpha
\quad\Longleftrightarrow\quad
\text{the action at coordinate }n\text{ of block }\alpha\text{ is specialization},
\]
equivalently,
\begin{align*}
   n\notin b_\alpha
\quad\Longleftrightarrow& \quad
\text{the action at coordinate }n\text{ of block }\alpha\text{ is forcing with }
T_{\omega\cdot\alpha+n} \\& \quad \text{and hence adding a cofinal branch}. 
\end{align*}
The construction gives the following dichotomy for every $\alpha\in\mathcal L_{\Sigma}$:
\[
b_\alpha=\Delta(u)\text{ for the unique real }u\text{ deliberately coded at }
\alpha,                                      \tag{2a}
\]
or else the block is a default block and
\[
b_\alpha=\omega.                             \tag{2b}
\]
Since $\omega\notin\operatorname{ran}(\Delta)$ and $\Delta$ is injective, the value of
$b_\alpha$, when it belongs to $\operatorname{ran}(\Delta)$, determines the coded real
uniquely.

Suppose first that $z$ is deliberately coded at the uniformization block
$\alpha\in\mathcal L_{\Sigma}$.  Let $R=R_\alpha$ be the real added by the
almost-disjoint coding step of this block.  Thus $R$ codes the three sequences
\[
\langle A_{\omega\cdot\alpha+n}:n\in\omega\rangle,
\qquad
\langle Y_{\omega\cdot\alpha+n}:n\in\omega\rangle,
\qquad
\langle T_{\omega\cdot\alpha+n}:n\in\omega\rangle,
\]
where $A_{\omega\cdot\alpha+n}$ is a specializing function if $n\in\Delta(z)$ and a
cofinal branch through $T_{\omega\cdot\alpha+n}$ if $n\notin\Delta(z)$.

Let $M$ be a countable suitable model with $R,z\in M$.  Decoding $R$ inside $M$ gives
initial segments
\[
\langle A^M_n:n\in\omega\rangle,
\qquad
\langle Y^M_n:n\in\omega\rangle,
\qquad
\langle U^M_n:n\in\omega\rangle,
\]
and a block index $\beta^M\in M$ which $M$ regards as an element of
$\mathcal L_{\Sigma}$.  By the localization part of the preliminary forcing
(the $K_3$-step in the FFDZ apparatus), for every $n\in\omega$ the model $M$ satisfies
that the tree decoded from the stationary-kill pattern at the block
$\omega\cdot\beta^M+n$ is precisely $U^M_n$.  In the notation of Definition~\ref{def:psi3},
\[
M\models
T^M_{\omega\cdot\beta^M,n}=U^M_n.              \tag{3}
\]
Moreover the object $A^M_n$ decoded from $R$ has, in $M$, exactly the kind prescribed by
$\Delta(z)$:
\[
\begin{array}{rcl}
n\in\Delta(z)     &\Longrightarrow& M\text{ sees that }A^M_n\text{ specializes }T^M_{\omega\cdot\beta^M,n},\\[2mm]
n\notin\Delta(z) &\Longrightarrow& M\text{ sees that }A^M_n\text{ is a cofinal branch through }T^M_{\omega\cdot\beta^M,n}.
\end{array}                                      \tag{4}
\]
Equations (3) and (4) verify the matrix in Definition~\ref{def:psi3}.  Since $M$ was
arbitrary, the real $R_\alpha$ witnesses $\PsiThree(z)$.

Conversely suppose that $\PsiThree(w)$ holds in $V[G]$, and let $R$ witness it.
Choose a sufficiently large regular $\Theta$ and a countable elementary submodel
\[
N\prec H_\Theta^{V[G]}
\]
with
\[
R,w,\Pp_{\omega_3},G,\mathcal L_{\Sigma},\Delta\in N.
\]
Let $\pi:N\to M$ be the transitive collapse.  Since $R$ and $w$ are reals, the
collapse fixes them.  The model $M$ is a countable suitable transitive model
containing $R$ and $w$, so Definition~\ref{def:psi3} supplies, inside $M$, an
index $\bar\beta\in\mathcal L_{\Sigma}^M$ witnessing the required
branch-versus-specialize pattern for $w$.  Let $\beta=\pi^{-1}(\bar\beta)$ and
let $\alpha$ be the corresponding external preliminary block index.

Inside $M$, for every $n\in\omega$, the tree decoded from the nonstationarity
pattern at the block $\omega\cdot\bar\beta+n$ is specialized exactly when
$n\in\Delta(w)$ and has a cofinal branch exactly when $n\notin\Delta(w)$.  By the
preliminary FFDZ localization and condensation analysis, this decoded tree is the
collapse of an initial segment of the genuine preliminary tree
$T_{\omega\cdot\alpha+n}$.  Consequently the pattern seen in $M$ is the collapse
of the actual final branch-versus-specialize pattern $b_\alpha$.  Hence
\[
\forall n\in\omega\quad
\bigl(n\in b_\alpha\Longleftrightarrow n\in\Delta(w)\bigr),
\]
and therefore
\[
b_\alpha=\Delta(w).                            \tag{5}
\]
The index $\alpha$ belongs to the uniformization-reserved class, because the
witnessing formula explicitly requires the block index to lie in
$\mathcal L_\Sigma$.  It cannot be a wellorder block, since the classes
$\mathcal L_{\mathrm{WO}}$ and $\mathcal L_\Sigma$ are disjoint.  It cannot be an
unused fresh block: by Lemma~\ref{lem:no-accidental-codes}, before a fresh block
is acted upon all trees in that block remain Suslin, so no
branch-versus-specialize pattern can be decoded there.  It cannot be a default
block either, because default blocks have pattern $\omega$, while
$\omega\notin\operatorname{ran}(\Delta)$.

Thus the action at block $\alpha$ was a genuine coding action.  By (2a), there is a
unique real $u$ deliberately coded at $\alpha$, and
\[
b_\alpha=\Delta(u).
\]
Together with (5) and the injectivity of $\Delta$, this gives $u=w$.  Hence $w$ was
deliberately placed into the uniformization-reserved block $\alpha$, as required.
\end{proof}

\end{document}
